% !TEX root = elm2022coldbursts.tex
% FICTIONAL TEST MATERIAL: this paper does not exist; it was written for the
% citation-audit fixture.
\documentclass{article}
\usepackage{amsmath}

\title{Cold-Weather Burst Rates in Iron Mains}
\author{H. Fixture \and J. Fixture}
\date{Journal of Fixture Studies, 2022}

\begin{document}
\maketitle

\noindent\textbf{FICTIONAL TEST MATERIAL: this paper does not exist; it was
written for the citation-audit fixture.} No author, utility, venue, or
measurement described here is real.

\begin{abstract}
Iron water mains break more often in winter. We combined eight years of burst
records from a fictional utility with soil-temperature measurements at the
depth of the mains. Burst rates in cast-iron and ductile-iron mains were 2.4
times higher on days when the soil temperature was below 2~$^\circ$C than on
days when it was above 2~$^\circ$C. We describe the threshold, its stability
across pipe classes, and its use for planning winter repair crews.
\end{abstract}

\section{Introduction}
\label{sec:introduction}

Cold weather stresses buried iron mains in two ways. Cold water contracts the
pipe and raises the axial stress at joints and corrosion pits, and freezing
soil shifts and loads the pipe. Utilities know that bursts cluster in winter,
but they lack a simple, measurable trigger for when the risk rises.

This paper relates burst rates to the soil temperature at the depth of the
mains, rather than to the air temperature, and identifies a threshold that
separates low-risk from high-risk days.

\section{Method}
\label{sec:method}

\subsection{Burst records}
\label{sec:method-records}

We used eight years of repair records from the fictional Fixture Heights
utility. A burst is a main break that required an unplanned repair. Only
cast-iron and ductile-iron mains were included. Burst counts were normalised by
the length of main in service on each day.

\subsection{Soil temperature}
\label{sec:method-soil}

Six soil-temperature stations recorded the temperature at 1.2~m, the typical
depth of the mains, every hour. Each burst was matched to the daily mean soil
temperature at the nearest station.

\subsection{Threshold analysis}
\label{sec:method-threshold}

We compared burst rates above and below candidate thresholds from
$-2$~$^\circ$C to 6~$^\circ$C in steps of 0.5~$^\circ$C and chose the threshold
that best separated the two groups.

\section{Results}
\label{sec:results}

\subsection{Burst-rate ratio at 2~$^\circ$C}
\label{sec:results-ratio}

The best threshold was a soil temperature of 2~$^\circ$C. Burst rates in cast-
and ductile-iron mains were 2.4 times higher when the soil temperature was
below 2~$^\circ$C than when it was above it. In other words, burst rates in
these iron mains more than doubled on the cold side of the threshold.

\subsection{Stability across pipe classes}
\label{sec:results-classes}

The ratio was similar for cast iron and ductile iron, and for small and large
diameters. It held in each of the eight winters of the record, although the
number of days below the threshold varied widely from winter to winter.

\subsection{Planning use}
\label{sec:results-planning}

Because soil temperature changes slowly, a forecast of the days below
2~$^\circ$C gives the utility several days of warning to schedule additional
repair crews.

\section{Limitations}
\label{sec:limitations}

This study analyses burst rates against soil temperature only.

This study does not examine chlorine or any other aspect of water quality: no
water samples were analysed. This study does not examine leak localization:
we used the repair records as given and did not consider how the bursts were
found or located. The analysis covers iron mains only; plastic mains respond to
cold differently and were excluded. The soil stations were few, so the
temperature matched to a burst may differ from the temperature at the exact
burst site.

\end{document}
